\section{Stability and the uniform arithmetic estimates}\label{paper:assembly}
\subsection{Stability of the original harmonic sum}
\label{sec:ell-stability}

The following comparison comes directly from the arithmetic sum. It
uses no Poisson bridge, saddle evaluation, or conjectured profile.
Write $c=\log3$ and $m_i=k-i+2$ as before, and fix an absolute constant
$C_{\rm M}$ for which
\[
 \left|\sum_{p\le t}\frac1p-\log\log t-B_1\right|
 \le C_{\rm M}\qquad(t\ge3).
\]
This form follows, for example, from the explicit estimates of
Rosser and Schoenfeld \cite{RosserSchoenfeld1962}; enlarging the absolute
constant covers the bounded range. Only this uniform coarse bound is used.

\begin{theorem}[Explicit log-parameter stability]\label{thm:ell-stability}
For $\ell,\ell'\in[c,\infty)^k$ put
\[
 \delta_i=\ell'_i-\ell_i,\quad
 D=\sum_{i=1}^k\left|\sum_{h=1}^i\delta_h\right|,
 \quad
 A=\sum_{i=1}^km_i|\delta_i|+(k+1)(D+2kC_{\rm M}).
\]
Then
\[
 e^{-A}K_k(\ell)\le K_k(\ell')\le e^A K_k(\ell).
\]
In particular, for each fixed $L<\infty$,
$K_k(\ell')\asymp_{k,L}K_k(\ell)$ uniformly whenever
$\|\ell'-\ell\|_1\le L$.
\end{theorem}
\begin{proof}
Let $y_i,y_i'$ be the associated cutoffs. Their definitions give
\[
 \log\log y_i'-\log\log y_i=\sum_{h\le i}\delta_h.
\]
Every prime that belongs to different blocks in the two configurations,
or belongs to only one of their prime ranges, lies in
\[
 \mathcal B=\bigcup_{i=1}^k
  \{p:\min(y_i,y_i')<p\le\max(y_i,y_i')\}.
\]
Indeed, outside this union the truth values of all inequalities
$p\le y_i$ and $p\le y_i'$ agree. Hence their least valid index,
if any, agrees. Mertens's estimate and positivity imply
\begin{equation}\label{eq:bad-prime-mass}
 \sum_{p\in\mathcal B}\frac1p\le D+2kC_{\rm M}.
\end{equation}

Retain only the common primes outside $\mathcal B$, with their common
block labels, and let $S_*$ be their unnormalized harmonic volume sum.
This number is positive, since the empty selection contributes
$(\log2)^k$. Write $S,S'$ for the two full unnormalized sums.
Inserting a selected block-$i$ prime gives the union of at most $m_i$
translates of the old union, one of them unchanged. Thus it increases
the volume by a factor between $1$ and $m_i\le k+1$.
Sum first over the common primes, and then over all selections of
the remaining primes, retaining their exact reciprocal weights. This
gives, for both $S_1=S$ and $S_1=S'$,
\[
 S_*\le S_1\le
 S_*\prod_{p\in\mathcal B}\left(1+\frac{k+1}{p}\right)
 \le S_*\exp\left((k+1)\sum_{p\in\mathcal B}\frac1p\right).
\]
It follows that $S'/S$ lies between the reciprocal and the value of
the last exponential; the exponent is not doubled. Finally the exact
ratio of the normalizing products is
$\exp(-\sum_i m_i\delta_i)$. Combine this identity with
\eqref{eq:bad-prime-mass} to obtain the assertion.
\end{proof}

\begin{corollary}[Contraction of any bounded collection of blocks]
\label{cor:bounded-blocks}
Fix $L\ge c$ and any set $I\subseteq\{1,\ldots,k\}$. If
$c\le\ell_i\le L$ for $i\in I$, replace those coordinates by $c$
and leave the others unchanged, obtaining $\widehat\ell$.
Then $K_k(\ell)\asymp_{k,L}K_k(\widehat\ell)$ uniformly over all
remaining coordinates. If $L$ is specified in terms of $k$ only, the
constants depend on $k$ only.
\end{corollary}
\begin{proof}
The changed vector remains admissible and its $\ell^1$ distance from
the original is at most $k(L-c)$. Apply
Theorem~\ref{thm:ell-stability}.
\end{proof}


\subsection{The lifted infimum and completion of the core estimate}
\begin{lemma}[Finite search range]\label{profile:finite-search}
Set
\[
 N=\frac{k(k+1)}2,\quad E=N+\frac{k-1}2+\frac32,
 \quad C=\frac{k(k+1)}2,\quad c=kH,\quad a=A_k-C,
\]
and
\[
 R=1+\max\left\{1,\max_i L_i,
          \frac{4E^2c}{a^2},
          \frac{2[A_k+(k-1)\log2+1]}a\right\}.
\]
The infimum in \eqref{profile:final} is unchanged if restricted to
$1\leq\lambda\leq R$. It is strictly positive and finite.
\end{lemma}
\begin{proof}
There are at most $N$ canonical edges. Since scores are at least one,
$Z_i\leq H$, and $2\chi\leq1$, each edge factor is at least
$(1+HL_i)^{-1}$. The bounds \eqref{profile:compact} give $0<d_j<1$.
Each proper count factor is therefore at least
$[2\sqrt{1+T_j}]^{-1}$. Finally
$\mathcal B\geq(1+V_k)^{-3/2}$. If $\lambda\geq\max_iL_i$,
every lifted length equals $\lambda$. The rate is at most
$\sum_jV_j=C\lambda$: in \eqref{profile:saddle}, $Z_i\geq1$
and $s_j\leq1$. Hence the objective is at least
\[
 2^{-(k-1)}e^{a\lambda}(1+c\lambda)^{-E}.
\]
The elementary inequality $\log(1+x)\leq\sqrt x$ gives
$a\lambda-E\log(1+c\lambda)\geq a\lambda/2$ when
$\lambda\geq4E^2c/a^2$. For $\lambda\geq R$ the last lower bound
exceeds $e^{A_k}$. At $\lambda=1$ the objective is at most $e^{A_k}$,
because $J\geq0$ and every factor in \eqref{profile:raw} is at most
one. This proves the range claim. The same explicit factor floors
on the bounded interval, and the elementary bound $J\leq kV_k$,
bound the infimum strictly away from zero. No continuity across an
owner merger is required for this argument.
\end{proof}


\begin{proposition}\label{profile:lift-removal}
The positive-length upper bound of Theorem~\ref{positive:upper},
the fixed-lift lower bound of Theorem~\ref{lower:lifted}, and
Theorem~\ref{thm:ell-stability} imply Theorem~\ref{main:uniform}.
\end{proposition}
\begin{proof}
Put $L_i^\lambda=\max(L_i,\lambda)$ and
$\delta_i=L_i^\lambda-L_i$, so that $0\leq\delta_i\leq\lambda$.
The exponent in Theorem~\ref{thm:ell-stability} satisfies
\[
 \sum_i h_i\delta_i\leq\frac{k(k+3)}2\lambda,
 \qquad
 \sum_{i=1}^k\left|\sum_{h\leq i}\delta_h\right|
 \leq\frac{k(k+1)}2\lambda.
\]
Hence it is at most $B_k\lambda+C_0(k)$, where
\[
 B_k=\frac{k(k^2+3k+4)}2,
 \qquad C_0(k)=2k(k+1)C_{\rm M},
 \qquad
 A_k-B_k=\frac{3k^3+9k^2+8k+4}{2}>0.
\]
For every $\lambda\geq1$, stability and the positive-length upper give
\[
 K_k(\ell)
 \leq e^{C_0(k)}C_U(k)e^{B_k\lambda}\Psi_k(L^\lambda)
 \leq e^{C_0(k)}C_U(k)e^{A_k\lambda}\Psi_k(L^\lambda).
\]
Taking the infimum proves $K_k(\ell)\ll_k\Phi_k(\ell)$.
For the reverse bound, choose just the value $\lambda=L_*(k)$
from Theorem~\ref{lower:lifted}:
\[
 \Phi_k(\ell)\leq e^{A_kL_*}\Psi_k(L^{L_*})
 \leq e^{A_kL_*}c_L(k)^{-1}K_k(\ell).
\]
All constants depend only on $k$. In particular, no comparison of
profiles at different equality faces is needed.
\end{proof}

\subsection{Transfer to the multiplication table}
\label{sec:table-transfer}

For ordered real parameters $3\le N_1\le\cdots\le N_{k+1}$, put
\[
 A_{k+1}(\mathbf N)
 =\#\{n_1\cdots n_{k+1}:n_i\in\mathbb N,\ n_i\le N_i\},
 \qquad P=\prod_{i=1}^{k+1}N_i.
\]
For ordered cutoffs $3\le t_1\le\cdots\le t_r$, write
$\mathcal K_r(\mathbf t)=K_r(\boldsymbol\ell(\mathbf t))$, where
\[
 \ell_i(\mathbf t)=\log\frac{3\log t_i}{\log t_{i-1}},
 \qquad t_0=3.
\]
The arithmetic normalization of this quantity is
\begin{equation}\label{eq:table-K-normalization}
 \mathcal K_r(\mathbf t)=
 \prod_{i=1}^r\left(\frac{\log t_i}{\log t_{i-1}}\right)^{-(r-i+2)}
 \sum_{\mathbf a\in\mathcal A_r(\mathbf t)}
 \frac{L^{(r+1)}(\mathbf a)}{a_1\cdots a_r}.
\end{equation}
Here $a_1$ is squarefree with prime factors at most $t_1$, and for
$i\ge2$ the prime factors of $a_i$ lie in $(t_{i-1},t_i]$.
In particular, the first prime block begins at $1$, not at $3$.
Write $\mathcal L^{(r+1)}(\mathbf a)$ for the union of boxes
\[
 \prod_{i=1}^r[\log(d_i/2),\log d_i),
 \qquad d_1\cdots d_i\mid a_1\cdots a_i\quad(1\le i\le r).
\]
Its volume is $L^{(r+1)}(\mathbf a)$.

\begin{theorem}[Full table transfer]\label{thm:full-table-transfer}
For every fixed integer $k\ge1$, define
\[
 Y_0=3,\qquad Y_i=\max\{3,N_i/2\}\quad(1\le i\le k),
 \qquad \lambda_i=\log\frac{3\log Y_i}{\log Y_{i-1}}.
\]
Then uniformly over $3\le N_1\le\cdots\le N_{k+1}$,
\begin{equation}\label{eq:table-K-transfer}
 A_{k+1}(\mathbf N)\asymp_k P K_k(\boldsymbol\lambda).
\end{equation}
Consequently, any independently proved comparison
$K_k(\boldsymbol\ell)\asymp_k\Phi_k(\boldsymbol\ell)$ on
$[\log3,\infty)^k$ gives
\[
 A_{k+1}(\mathbf N)\asymp_k P\Phi_k(\boldsymbol\lambda)
\]
on this entire table region. This implication requires no stability
assumption on $\Phi_k$.
\end{theorem}

\begin{proof}
We use two published reductions, whose hypotheses we record explicitly.
Koukoulopoulos's Theorem~1.1 \cite{Kouk} gives, for every $r\ge1$ and every
$3\le M_1\le\cdots\le M_{r+1}$,
\[
 A_{r+1}(\mathbf M)\asymp_r
 H^{(r+1)}\left(\prod_{i=1}^{r+1}M_i,
                   (M_1/2,\ldots,M_r/2),(M_1,\ldots,M_r)\right).
\]
His Theorem~1.7 \cite{Kouk} gives $H^{(r+1)}(x,\mathbf t,2\mathbf t)/x
\asymp_r\mathcal K_r(\mathbf t)$ provided
$3\le t_1\le\cdots\le t_r$ and
$x\ge2^r t_1\cdots t_rt_r$.
Neither reduction uses a proposed evaluation of $K_r$.

If $M_1\ge6$, then $t_i=M_i/2\ge3$ and
\[
 2^r t_1\cdots t_rt_r
 =\left(\prod_{i=1}^rM_i\right)\frac{M_r}{2}
 \le\frac12\prod_{i=1}^{r+1}M_i.
\]
Thus the two reductions apply directly. Choose constants $a_r,b_r>0$
depending only on $r$ such that
\begin{equation}\label{eq:table-large-axis}
 a_r\mathcal K_r(M_1/2,\ldots,M_r/2)
 \le\frac{A_{r+1}(\mathbf M)}{\prod_iM_i}
 \le b_r\mathcal K_r(M_1/2,\ldots,M_r/2)
 \quad(M_1\ge6).
\end{equation}

Let $b=\#\{1\le i\le k:N_i<6\}$ and $m=k-b$.
By ordering these are the first $b$ indices. Set
\[
 P_{\rm tail}=\prod_{i=b+1}^{k+1}N_i,\qquad
 A_{\rm tail}=A_{m+1}(N_{b+1},\ldots,N_{k+1}),
\]
with $A_1(N)=\lfloor N\rfloor$.
The full table contains the tail table, by taking its first $b$
factors equal to $1$. Conversely it is a union of at most
$\prod_{i=1}^b\lfloor N_i\rfloor$ integer dilates of the tail table.
Since $\prod_{i=1}^bN_i<6^b$ when $b>0$, we have in all cases
\begin{equation}\label{eq:table-delete-axes}
 6^{-b}\frac{A_{\rm tail}}{P_{\rm tail}}
 \le\frac{A_{k+1}(\mathbf N)}P
 \le\frac{A_{\rm tail}}{P_{\rm tail}}.
\end{equation}

Suppose $m\ge1$ and put $z_j=N_{b+j}/2$ for $1\le j\le m$.
Then $z_1\ge3$, and \eqref{eq:table-large-axis} applies in dimension
$m$. We claim, on writing $h=\log2$ and $D=\log12$, that
\begin{equation}\label{eq:table-insert-three}
 h^b\mathcal K_m(\mathbf z)
 \le\mathcal K_k(\underbrace{3,\ldots,3}_{b},\mathbf z)
 \le D^b\mathcal K_m(\mathbf z).
\end{equation}
For $b=0$ this is equality. For $b\ge1$, every tuple on the middle
side satisfies $a_1\mid6$ and $a_2=\cdots=a_b=1$.
The map
\[
 (a_1,1,\ldots,1,a_{b+1},\ldots,a_k)
 \longmapsto(a_1a_{b+1},a_{b+2},\ldots,a_k)
\]
is a bijection onto $\mathcal A_m(\mathbf z)$, and preserves the
weight $1/(a_1\cdots a_k)$: the factors at primes $2,3$ uniquely
recover $a_1$. Let $\mathbf a'$ denote its image. Setting the first
$b$ divisor coordinates equal to $1$ gives the inclusion
\[
 [-h,0)^b\times\mathcal L^{(m+1)}(\mathbf a')
 \subseteq\mathcal L^{(k+1)}(\mathbf a).
\]
Conversely each of the first $b$ divisor coordinates is at most $6$,
and the projection onto the last $m$ coordinates belongs to
$\mathcal L^{(m+1)}(\mathbf a')$. Hence
\[
 \mathcal L^{(k+1)}(\mathbf a)
 \subseteq[-h,\log6)^b\times\mathcal L^{(m+1)}(\mathbf a').
\]
Taking volumes gives the factors $h^b,D^b$. The first $b$
normalizing log ratios are $1$; the remaining ratios and their
exponents agree exactly with those in dimension $m$.
Summing proves \eqref{eq:table-insert-three}.
Together with \eqref{eq:table-large-axis} and
\eqref{eq:table-delete-axes}, this yields
\[
 6^{-b}a_mD^{-b}\mathcal K_k(\mathbf Y)
 \le\frac{A_{k+1}(\mathbf N)}P
 \le b_mh^{-b}\mathcal K_k(\mathbf Y).
\]

It remains to handle $m=0$. The ratio
$A_{\rm tail}/P_{\rm tail}=\lfloor N_{k+1}\rfloor/N_{k+1}$ lies in
$[2/3,1]$. At $\mathbf Y=(3,\ldots,3)$ the only tuples are
$(a_1,1,\ldots,1)$ with $a_1\mid6$, and their reciprocal weights
sum to $2$. Every such union has volume between $h^k$ and $D^k$.
Consequently
\[
 2h^k\le\mathcal K_k(3,\ldots,3)\le2D^k,
\]
and \eqref{eq:table-delete-axes} gives the explicit bounds
\[
 \frac{\mathcal K_k(\mathbf Y)}{3(6D)^k}
 \le\frac{A_{k+1}(\mathbf N)}P
 \le\frac{\mathcal K_k(\mathbf Y)}{2h^k}.
\]
Taking the minimum and maximum over the finitely many possible
$b\in\{0,\ldots,k\}$ proves \eqref{eq:table-K-transfer} with
constants depending only on $k$.
\end{proof}


\begin{proof}[Proof of Theorems~\ref{paper:main-H} and~\ref{paper:main-table}]
Theorem~\ref{main:uniform} and the arithmetic comparison
\eqref{main:admissible} give the asserted estimate for
$H^{(k+1)}$ throughout the stated region. For the table, apply
Theorem~\ref{thm:full-table-transfer} at the original axes and then
Theorem~\ref{main:uniform} at $\ell(\mathbf Y)$. Both comparisons have
constants depending only on the fixed dimension. Adjacent equal cutoffs,
bounded axes, and arbitrarily small strict saddle gaps are included in
these same arguments.
\end{proof}
